\pagebreak

\chapter{Immediate Statements}\label{chap:immed}

    Immediate statements generate no target code execution threads.

    An immediate variable is one which has been declared with no mode, or
    explicitly with immediate mode.


    \section{immediate mode assignment}\label{sec:im_ass}

        \index{assignment!immediate mode}
        Immediate assignment is assignment to an immediate variable.
        An immediate assignment statement has the forms -
        \begin{lstlisting}[gobble=8, language=threepl]
               variable = expression;          // simple assignment
               variable += expression;         // additive assignment
               variable -= expression;         // subtractive assignment
               variable *= expression;         // multiplicative assignment
               variable /= expression;         // divisional assignment
               variable %= expression;         // remainder assignment
               variable &&= expression;        // logical AND assignment
               variable ||= expression;        // logical OR assignment
               variable ^^= expression;        // logical XOR assignment
               variable &= expression;         // bitwise AND assignment
               variable |= expression;         // bitwise OR assignment
               variable ^= expression;         // bitwise XOR assignment
               variable <<= expression;        // left shift assignment
               variable >>= expression;        // right shift assignment
               variable++;                     // increment statement
               variable--;                     // decrement statement
        \end{lstlisting}
        Where a variable is a structure or an array, members can
        be assigned by -
        \begin{lstlisting}[gobble=8, language=threepl]
               variable.field = expression;
               variable[subscript] = expression;
               variable.field++;
               variable.field--;
               variable[subscript]++;
               variable[subscript]--;
        \end{lstlisting}
        
        For simple assignment whole structs or arrays can be assigned by
        \begin{lstlisting}[gobble=8, language=threepl]
               variable = variable;
               variable = {.......};   // compound value
        \end{lstlisting}
        
        \index{string/type conversion}
        For simple assignment if the LHS is type "str" and the RHS is
        type "type" the RHS will be converted to string representation of the
        type (the same as invoking unbuilt function {\bf typestring()}).
        If the LHS is type "type" and the RHS is type "str" the string will
        be converted to a type.

        For the other assignment types the RHS type must be primitive. Types
        allowed are -

        \begin{tabular}{| l | l | l |}
        \hline
                  & LHS & RHS  \\
        \hline
        \hline
        \verb'+=' & int, uint & int, uint  \\
                  & float & float  \\
                  & float & int, uint  \\
                  & str & str, int, uint, float  \\
        \hline
        \verb'-=' & int, uint & int, uint \\
                  & float & float  \\
                  & float & int, uint  \\
        \hline
        \verb'*=' & int, uint & int, uint \\
                  & float & float  \\
                  & float & int, uint  \\
        \hline
        \verb'/=' & int, uint & int, uint \\
                  & float & float  \\
                  & float & int, uint  \\
        \hline
        \verb'%=' & int, uint & int, uint \\
                  & float & int, uint  \\
        \hline
        \verb'&&=' & log & log \\
        \hline
        \verb'||=' & log & log \\
        \hline
        \verb'^^=' & log & log \\
        \hline
        \verb'&=' & int, uint, bits & int, uint, bits \\
        \hline
        \verb'|=' & int, uint, bits & int, uint, bits \\
        \hline
        \verb'^=' & int, uint, bits & int, uint, bits \\
        \hline
        \verb'<<=' & int, uint, bits & int, uint \\
        \hline
        \verb'>>=' & int, uint, bits & int, uint \\
        \hline
        \end{tabular}

        Note that \verb'+=' can be used for string concatenation where the
        RHS can be one of several types.
        
        The logical assignments \verb'&&=' and  \verb'||=' do not perform
        short-circuit evaluation unlike the equivalent binary operators.

        RHS expressions may contain only immediate variables and constants
        and related operators and function calls returning immediate values.
        
        LHS variable references may have side effects, for example incremented
        or decremented subscripts or subscripts returned by calls to functions
        within which enduring state is changed.
        
        As a special case a value mode expression whose value is a constant
        may be assigned to an immediate mode variable. A value mode expression
        may be tested by calling the inbuilt function {\bf istargconst()}
        \autorefph{sec:ifistargconst}. These features may be used where
        either a constant of a variable is passed as an argument to a parameter
        and immediate conditional code needs to distinguish which is the case.


    \section{value mode assignment}\label{sec:im_valass}

        \index{assignment!value mode}
        Assignment to a value mode variable.
        is an immediate assignment, it is not executed in the target. It may
        however create target logic.
        A value assignment statement has the form -
        \begin{lstlisting}[gobble=8, language=threepl]
               variable := expression;
        \end{lstlisting}
        
        Whole structs or arrays can be assigned by
        \begin{lstlisting}[gobble=8, language=threepl]
               variable := variable;
        \end{lstlisting}

        RHS expressions may contain constants, both immediate and target
        variables and function calls and operators.
        
        A value variable is a target variable which represents the value of a
        target expression. Constants will be converted to equivalent signals in
        the target logic.
        
        A value variable can also be assigned an initial value when created
        (see {\bf value()} \autorefph{sec:ipvalue})
        
        Whether the assignment is of a primitive or a compound type the widths
        of individual words may differ as long as the primitive types match.
        Bits and unsigned integers will be padded or truncated as needed.
        Signed integers will be truncated or sign extended. Signed integers
        may not be assigned to unsigned integers. Unsigned integers may be
        assigned to signed integers, but since size adjustments are made
        by truncation or zero padding the sign of the result may be incorrect
        unless the LHS is at least one bit larger than the RHS.
        
        A value mode variable may be created as type "empty", in which case it
        can be assigned any type. Following assignment the variable will retain
        the type of the RHS value assigned to it. Any further assignments to
        that variable will be subject to type checking as normal. A value
        variable can be cleared and its type changed to "empty" by calling
        inbuilt procedure {\bf emptyvalue} \autorefph{sec:ipemptyvalue}.
        
        A value variable is usually combinatorial in that it does not itself have
        intrinsic state in the target device but may explicitly create state
        devices. For example in the following
        \begin{lstlisting}[gobble=8, language=threepl]
               value(v, "log");
               
               v := a && b || c;
        \end{lstlisting}
        variable {\em v} represents the value of the logical expression and as
        such is combinatorial.
        
        But in the following example
        \begin{lstlisting}[gobble=8, language=threepl]
               value(v, "log");
               
               v := a;
        \end{lstlisting}
        {\bf v} is effectively a pointer to variable {\em a} (though not one that
        can be used with the indirect operator). The state, or storage, is in {\em
        a}, not in {\em v}. If {\em a} changes, the value of {\em v} will also
        change since it simply reflects the value of {\em a}.
        
        Value variables are used mainly as 'pointers' to combinatorial
        expressions that are to be shared. 3PL will not attempt to find
        common subexpressions and instantiate these once, indeed that can be
        counter-productive for speed. Value variables as parameters are also
        a way of passing a variable or expression to or from a module,
        procedure or function.
            
        When a value variable is a compound type each element of the
        type is separate. Thus a value array can be linked to
        quite distinct sources, e.g.
        \begin{lstlisting}[gobble=8, language=threepl]
               value(a, "[3]uint:8");
               static(b, "uint:8");
               queue(c, "uint:8");
               value(d, "[2]uint:8");

               d[1] := a + 2;
               a[0] := b;
               a[1] := c;
               a[2] := 2 * d[1];
        \end{lstlisting}
        The value array {\em a} could now be assigned in total to
        another value array, whole or in part -
        \begin{lstlisting}[gobble=8, language=threepl]
               value(e, "[5]uint:8");

               e[1..3] := a;
        \end{lstlisting}
        Each component of the left hand side value variable array is
        assigned the value of the corresponding right hand side value
        variable.
        
        Value variables can be re-assigned at any time. When a value variable
        is referenced its current value is used. If a value variable is
        referenced prior to its first assignment, a temporary link will be put
        in place and this will be replaced by the first assigned value after
        which the value variable behaves as described above. This
        pre-assignment usage is provided to avoid problems associated with
        cycles of dependence which occur frequently.
        
        If the value of a value variable is used prior to assignment the
        value will be given the current clock domain and the value assigned
        later must have the same clock domain.
        
        Where a value variable is assigned a value following evaluation of that
        variable, i.e. it has been used prior to assignment, the value assigned
        must not contain queue reads.
        
        Where a value variable is assigned a value following evaluation of that
        variable, i.e. it has been used prior to assignment, the value assigned
        must not contain the function {\bf cmemoryread()}
        \autorefph{sec:ifcmemoryread}. \footnote{This restriction is imposed
        because of the need to link the execution signal for a statement
        containing a {\em cramread()} (or {\em cmemoryead()}) to the memory
        variable itself. The execution signal is used to multiplex the read
        addresses where multiple reads are involved. Post assignment to a value
        variable occurs after that execution signal was available. This
        limitation may be removed in the future}.
        
        A value must not be assigned an expression containing {\bf
        cmemoryread()} \autorefph{sec:ifcmemoryread} if the value is never used
        as this generates a read access to the memory variable which is never
        connected to appropriate control logic. This restriction may be removed
        in future.
        
        There is a special case of value assignment which allows a clock variable
        to be used as a logic value. A clock variable can be assigned to a value
        mode variable of type {\bf log}.
        
        A value variable as a whole has a no clock domain, but the components of a
        value will each have a clock domain. That clock domain will be determined
        when the variable is evaluated by examining the clock domains of all
        ultimate source variables in expressions, i.e. variables that have state
        such as static mode and queue mode. The components, i.e. array members or
        structure fields, need not have the same clock domain, though that would
        be unusual. When a value mode variable is evaluated the value will have
        the clock domain of the component. If a number of components are accessed
        in one evaluation, for example a sub-array, they must all have the same
        clock domain.
        
        A value variable may have the clock domains of its components set by
        passing the {\bf domain} attribute in the variable creation procedure call
        or by setting it by calling inbuilt procedure {\bf attributes()}. This will
        establish the clock domains of any components that thus far have no
        clock domain, all the way back to variables with state such as static mode
        or queue mode variables. Where such variables already have a clock domain
        the domain will be compared and if different a fatal error will result.
        
        For a discussion of clock domains see \autorefph{subsec:clockdoms}.
        



    \section{priority mode assignment}\label{sec:priass}
    
        \index{assignment!priority mode}
        Priority mode variables are always arrays of type {\bf log}. An
        assignment to a member, or members, of a priority variable array forms
        an input to the priority encoder where subscript 0 has the highest
        priority.
        
        Assignment of priority mode variable array members has the same syntax
        as for value variables -
        \begin{lstlisting}[gobble=8, language=threepl]
               variable[i] := expression;
        \end{lstlisting}
        
        The assigned expression, which is an input, must not contain queue
        reads, but it may contain queue examines and queue availability tests.
        
        A priority variable array member may be assigned more than once but the
        default behaviour is that the final value assigned prior to termination
        of immediate execution prevails, all other previous assignments being
        lost. This can be changed by setting attribute "multiplein" to true in
        which case all assignments will be collected and ORed together on
        completion of immediate execution to form the input to that priority
        array member.
            
        A priority variable inherits the current clock domain. Values assigned
        to it must either have the same clock domain. If the value to be
        assigned has a different clock domain it must be converted by one of
        the inbuilt functions {\bf sample()} \autorefph{sec:ifsample} or {\bf
        accept()} \autorefph{sec:ifaccept}.

        A priority variable has a single clock domain for reading and writing.
        When a priority variable is created using the inbuilt procedure {\bf
        priority()} \autorefph{sec:ippriority} it has no associated clock domain.
        The clock domain is established by
        \blist
        \item   the first assignment to the variable
        \item   the first use of the value of the variable
        \blend
        
        {\bf Note that evaluation of the variable in any context, for example
        in inbuilt or library modules, procedures or functions, will establish
        a clock domain if the variable does not yet have a clock domain.}
        This will always be the current clock domain (unless the code has
        been explicitly written to establish the clock domain of its argument
        to something other than the current clock domain, in which case
        this will presumably be well documented).
        
        When a priority variable is assigned, the clock domain of the assigned
        value must be the same as the current clock domain or null.
        
        When a priority variable is assigned, If it has no clock domain the
        current clock domain will be established as the clock domain of the
        variable. If it has a clock domain it must be the same as the current
        clock domain.
        
        When the value of a priority variable is used that value will have the
        clock domain of the variable.
        
        Note that control structues {\bf waitpri} and {\bf waitprilock} use
        priority variables for mutual exclusion and provide a more convenient
        interface than use of priority variables directly
        \autorefph{sec:waitpri}.
        
        



    \section{output mode unconditional assignment}\label{sec:outputiass}
    
        \index{assignment!unconditional!output mode}
        An output variable represents a connection off-chip via output buffers and
        pads. Unconditional assignment of output mode variables is used where
        3-state output is not required.

        Unconditional assignment has the same syntax as for value mode variables -
        \begin{lstlisting}[gobble=8, language=threepl]
               variable := expression;
        \end{lstlisting}
        
        Any component of an output mode variable may only be assigned once,
        either by default via the 3rd argument to inbuit procedure output() or
        explicitly by a conditional or unconditional assignment statement.
        Components such as structure fields or array members may be assigned
        in separate statements and may be a mixture of conditional and
        unconditional.
                
        In the Xilinx implementation if the source of the output is a static
        variable then the associated flip-flops can be located in the IOBs. This
        can be done by a suitable flag to the Xilinx tools but is best done by
        explicitly setting the IOB attribute to true on the static variable since
        the 3PL netlist optimiser is then aware of this. If the static varible
        flip-flops go to other destinations as well as the output buffers then they
        will be replicated by the 3PL netlist optimiser, the replicants being
        sourced in parallel and located in IOBs while the originals stay in CLBs.
        The advantage of this is that the propagation time from flip-flop to pad
        will be the minimum achievable (controlled also by slew rate and drive
        current attributes). If this is not the case then a timing constraint from
        the clock domain to the output pads may need to be provided.
        
        If a source for an output mode variable is specified when it is created
        (3rd argument to {\bf output()}, see \autorefph{sec:ipoutput}) that
        is equivalent to an initial unconditional assignment of that source
        to the output mode variable.
        
        The assigned expression must not contain queue reads, but it may contain
        queue examines and queue availability tests. The LHS may have array
        subscripts, subscript ranges or struct fields. If the value to be assigned
        has a different clock domain it must be converted by one of the inbuilt
        functions {\bf sample()} \autorefph{sec:ifsample} or {\bf accept()}
        \autorefph{sec:ifaccept}.
          
        An output variable has a write clock domain but no read clock domain.
        When an output variable is created using the inbuilt procedure {\bf
        output()} \autorefph{sec:ipoutput} it has no associated clock domain. Its
        clock domain is established by the first assignment to the variable (or
        component of the variable if the type is compound) and is the current
        clock domain. The RHS exression must also have the current clock domain
        as its output (read) clock domain. The RHS expressions may be immediate
        or target mode. 
        
        An output mode variable must have at least a {\bf loc} attribute,
        providing a location for each pad.
        
        



    \section{clock mode assignment}\label{sec:clockass}
    
        \index{assignment!clock mode}
        Assignment of clock mode variables has the same syntax as for value
        variables -
        \begin{lstlisting}[gobble=8, language=threepl]
               variable := variable;
        \end{lstlisting}
        but the RHS can only be a clock variable or a function returning
        a clock variable. An assignment of null is not allowed.
            
        The LHS clock variable now becomes an implicit indirect reference to
        the RHS clock variable. The two variables are now equivalent in all
        contexts (but the netlist identifier is the RHS variable).
        
        These assignments may be chained.
            
        Where clock variables are assigned and are also used to set the clock
        domain, note that logic using that clock will be connected to the
        clock source determined by the value the clock variable has when
        immediate execution terminates. Changing the value of a clock variable
        by assignment, even if accompanied by passing it to {\bf useclock()},
        will not result in associated code having different clock domains -
        the code will all be driven by the clock source given by the final
        value of the clock variable.
        



    \section{the {\bf if} control structure}\label{sec:im_if_struct}

        \index{statement!if}
        The {\bf if} has the format -
        \begin{lstlisting}[gobble=8, language=threepl]
               if (expression)
                   body
        \end{lstlisting}
        which executes the body if the test expression is true.

        An alternative form is 
        \begin{lstlisting}[gobble=8, language=threepl]
               if (expression)
                   true_body
               else
                   false_body
        \end{lstlisting}
        which provides an additional body to be executed when
        the test expression is false.

        The expression must be immediate and must evaluate to a logical
        result.

        In the first form the body will be executed if the
        expression is true. In the second form one of the code
        bodies will be executed as selected.
        
        A body may be a single statement or a statement block in braces ({\em
        \{\dots\}}).
        
        Any target statements generated by execution of one of the bodies will
        replace this {\em if} statement in the target code. If there are
        multiple target statements generated then this {\em if} must be
        enclosed in a {\em seq}, {\em par} or {\em sync} block to control
        their execution.



    \section{the {\bf while} control structure}\label{sec:im_while_struct}

        \index{statement!while}
        The {\bf while} control structure iterates while
        a test expression is true. The format is -
        \begin{lstlisting}[gobble=8, language=threepl]
               while (expression)
                   body
        \end{lstlisting}
        which executes the body while the test expression
        remains true.

        The expression must be immediate and must evaluate to a logical
        result.
        
        The body may be a single statement or a statement block in braces ({\em
        \{\dots\}}).
        
        Any target statements generated by execution of the body will
        accumulate as a list of statements to replace this {\em while}
        statement in the target code. If there are multiple target statements
        generated then this {\em while} must be enclosed in a {\em seq}, {\em
        par} or {\em sync} block to control their execution.



    \section{the {\bf do while} control structure}\label{sec:im_dowhile_struct}

        \index{statement!do while}
        The {\bf do while} control structure iterates while
        a test expression is true. The format is -
        \begin{lstlisting}[gobble=8, language=threepl]
               do {
                   body
               } while (expression);
        \end{lstlisting}
        which executes the body while the test expression
        remains true. This differs from the {\bf while} in
        that the body is executed at least once.

        The expression must be immediate and must evaluate to a logical
        result.
        
        Any target statements generated by execution of the body will
        accumulate as a list of statements to replace this {\em do while}
        statement in the target code. If there are multiple target statements
        generated then this {\em do while} must be enclosed in a {\em seq},
        {\em par} or {\em sync} block to control their execution.



    \section{the {\bf for} control structure}\label{sec:im_for_struct}

        \index{statement!for}
        The {\bf for} control structure iterates while a test expression is
        true and provides initialisation and iteration control. The format
        is -
        \begin{lstlisting}[gobble=8, language=threepl]
               for (initialisation ; expression ; iteration)
                   body
        \end{lstlisting}
        The initialisation statement must be a single immediate assignment
        statement or an inbuilt variable creation procedure call. If the latter
        the resulting variable will be in a local scope containing the loop.

        The expression must be immediate and must evaluate to a logical
        result.
        
        The iteration statement must be a single immediate assignment statement
        or increment or decrement statement as the C ',' (comma) operator is not
        implemented in 3PL.

        The body may be a single statement or a statement block in braces
        ({\em \{\dots\}}).

        Any target statements generated by execution of the body will
        accumulate as a list of statements to replace this {\em for} statement
        in the target code. If there are multiple target statements generated
        then this {\em for} must be enclosed in a {\em seq}, {\em par} or {\em
        sync} block to control their execution.



    \section{the {\bf switch} control structure}\label{sec:im_switch_struct}

        \index{statement!switch}
        The {\bf switch} statement selects a code block
        by testing the value of an expression against a
        number of cases. The format is -
        \begin{lstlisting}[gobble=8, language=threepl]
               switch (expression) {
               case expression_or_list1:
                   body1
               case expression_or_list2:
                   body2
               default:
                   bodyd
               }
        \end{lstlisting}
        each {\em expression\_or\_list} may be a single expression or a
        comma-separated list of expressions enclosed in braces. The value
        derived from {\em expression} is compared with the values in each of
        the case statements and equality results in the associated code body
        being executed. {\bf expression} may be any immediate primitive type,
        including float\footnote{floating point equality comparisons require
        caution as rounding errors in computed values often result in values
        differing from the exact}. The case values must be of the same type as
        the test expression, but they need not evaluate to constants. The
        case values, whether constant or not, need not be unique. If they are
        repeated the first one encountered (from the top) will be used.
        
        {\bf Note that unlike languages such as C or Java the code associated
        with a case value does not flow on to the next case, i.e. break
        statements are ignored by the switch and may be used within a
        switch case code body to exit from an enclosing loop.}
        
        Any target statements generated by execution of a body will replace
        this {\em switch} statement in the target code. If there are multiple
        target statements generated then this {\em switch} must be enclosed in
        a {\em seq}, {\em par} or {\em sync} block to control their execution.



    \section{the {\bf break} statement}\label{sec:im_break_statement}

        \index{statement!break}
        A {\bf break} statement can only occur within immediate {\bf while},
        {\bf do while} or {\bf for} loops. It switches immediate execution to
        the end of the enclosing iterative statement, exiting that statement.



    \section{the {\bf continue} statement}\label{sec:im_continue_statement}

        \index{statement!continue}
        A {\bf continue} statement can only occur within immediate {\bf while},
        {\bf do while} or {\bf for} loops. It switches immediate execution to
        the end of the enclosing iterative statement to conditionally execute
        the next iteration.



    \section{the {\bf return} statement}\label{sec:im_return_statement}

        \index{statement!return}
        A {\bf return} statement returns from a call to a module, procedure
        or function. It may also be used to terminate execution of a file module
        (which is executing via an implicit call) or to terminate execution
        of an inline module. On returning from the initial file module any
        trailer modules will then be executed.

        For a module or procedure the syntax is -
        \begin{lstlisting}[gobble=8, language=threepl]
               return;
        \end{lstlisting}
        For a function the syntax is -
        \begin{lstlisting}[gobble=8, language=threepl]
               return(v);
        \end{lstlisting}
        where the argument is an expression giving the returned value of the
        function.
        
        A module or procedure will of course return to the point
        at which it was called if code execution simply proceeds past the
        last statement in its code block. Functions must use {\bf return}
        to return a value.
        
        A {\bf return} statement cannot appear in the following target
        control structures -
        \blist
        \item   branch
        \item   par
        \item   seq
        \item   sync
        \item   sync when
        \item   seq do while
        \item   par do while
        \item   sync do while
        \item   seq while
        \item   par while
        \item   sync while
        \item   unless
        \item   waitpri
        \item   when
        \blend


